\chapter{Introducción}
\label{ch:introduccion}

\section{Contexto}

Los sistemas de información contemporáneos deben atender simultáneamente necesidades
que antes solían resolverse en plataformas separadas. Una misma aplicación puede
registrar operaciones de negocio, reaccionar a eventos en tiempo real, conservar un
historial auditable, aislar información de varias organizaciones y producir datos para
analítica o inteligencia artificial. Esta convergencia ha aumentado el valor de las
bases de datos NoSQL, cuya flexibilidad permite representar estructuras heterogéneas y
distribuir grandes volúmenes de información. Sistemas influyentes como Bigtable
demostraron que un modelo distribuido orientado a columnas puede soportar cargas
muy diferentes cuando el diseño físico se alinea con los patrones de acceso
\parencite{chang2006bigtable}. Dynamo, por su parte, expuso de forma explícita los
compromisos entre disponibilidad, particionamiento, replicación y consistencia en un
almacenamiento distribuido orientado a clave \parencite{decandia2007dynamo}.

Esa flexibilidad no elimina la necesidad de modelar; la desplaza. En sistemas NoSQL,
la aplicación debe decidir qué datos se almacenan juntos, cuáles se duplican, qué
índices se mantienen y cómo se reparan las representaciones derivadas. El esquema puede
ser flexible en la base, pero sigue existiendo en el código, en las rutas, en las reglas
de acceso y en las expectativas de cada consumidor.

La dificultad aumenta en aplicaciones \emph{multiinquilino}, es decir, aquellas en las
que una misma plataforma sirve a varias organizaciones o contextos lógicos. La
consolidación puede reducir costos, pero introduce requisitos estrictos de aislamiento,
configuración, rendimiento y recuperación. La literatura ha mostrado que no existe un
único diseño físico que sea óptimo para todos los organizaciones y cargas; las alternativas
intercambian escalabilidad, espacio, flexibilidad y complejidad operativa
\parencite{aulbach2009flexibleschemas,yaish2014adaptiveschema,rico2016multitarget}.

Al mismo tiempo, los datos operacionales han dejado de ser únicamente un registro del
pasado. Se utilizan para detectar anomalías, descubrir secuencias, construir modelos
predictivos, recuperar evidencias y asistir decisiones humanas. Los sistemas de
aprendizaje automático necesitan características consistentes entre entrenamiento e
inferencia, versiones reproducibles y trazabilidad hacia las fuentes. Los almacenes de
características surgieron precisamente para organizar este flujo y reducir la
fragmentación entre preparación, entrenamiento y servicio de modelos
\parencite{orr2021featurestores,li2023featurestore}. Sin embargo, una capa analítica no
puede corregir por sí sola un origen operacional ambiguo o lleno de duplicaciones sin
propietario.

Este trabajo parte de una observación práctica: los sistemas NoSQL de alcance amplio
necesitan distinguir con claridad qué representación conserva el estado autorizado del
negocio y cuáles representaciones existen para acelerar consultas, integrar módulos o
alimentar procesos analíticos. A partir de esa separación se formula \COBRIA{}, una
arquitectura que organiza entidades canónicas, mapeos físicos, repositorios, servicios,
casos de uso y proyecciones reconstruibles dentro de una misma disciplina de diseño.

\section{Problema de investigación}

En un modelo desnormalizado, una modificación lógica puede afectar varias ubicaciones.
Por ejemplo, actualizar el estado de una entidad puede requerir modificar su registro
principal, retirar una referencia de un índice, añadirla a otro y actualizar un resumen.
Si estas ubicaciones se tratan como fuentes equivalentes, el sistema pierde una base
clara para decidir cuál valor es correcto cuando aparece una discrepancia. Si, por el
contrario, se reconoce una fuente canónica y las demás se definen como funciones
derivadas, la discrepancia puede detectarse y repararse.

La vista materializada es un antecedente directo de esta idea: una representación
derivada puede mejorar el acceso, pero debe mantenerse frente a cambios en los datos
base. La investigación clásica sobre mantenimiento incremental estudió cómo calcular
los cambios de una vista sin reconstruirla completamente
\parencite{gupta1993incremental,gupta1995materialized}. En aplicaciones NoSQL, el mismo
problema reaparece en índices manuales, resúmenes, relaciones inversas, contadores,
colas de salida y conjuntos de datos analíticos, aunque no siempre se describa con esa
terminología.

Se identifican seis manifestaciones principales del problema:

\begin{enumerate}[label=\textbf{P\arabic*.},leftmargin=*]
  \item \textbf{Ambigüedad de autoridad.} Varias copias contienen atributos similares,
  pero no se declara cuál es la fuente autorizada.

  \item \textbf{Consultas no acotadas.} Ante la ausencia de índices apropiados, el
  cliente recupera colecciones grandes y filtra localmente, aumentando transferencia,
  memoria, latencia y exposición innecesaria.

  \item \textbf{Derivaciones frágiles.} Los índices se actualizan mediante lógica
  dispersa y carecen de versión, propietario, verificador o procedimiento de
  reconstrucción.

  \item \textbf{Acoplamiento semántico--físico.} Nombres, claves compactas, rutas y
  reglas se duplican en múltiples capas; una evolución del esquema puede dejar
  consumidores incompatibles.

  \item \textbf{Aislamiento incompleto.} El organización se aplica como filtro de interfaz,
  pero no como dimensión obligatoria del modelo, la consulta y la autorización.

  \item \textbf{Preparación analítica no reproducible.} Los conjuntos de datos y características
  se crean con exportaciones o scripts independientes, sin linaje suficiente hacia los
  datos operacionales y sin corrección temporal verificable.
\end{enumerate}

Estas manifestaciones están relacionadas. Una consulta amplia suele motivar la creación
de un índice; el índice introduce duplicación; la duplicación exige mantenimiento; el
mantenimiento necesita autoridad y trazabilidad; y la trazabilidad se vuelve todavía
más importante cuando el resultado alimenta un modelo de IA. El problema de
investigación se formula, por tanto, de la siguiente manera:

\begin{quote}
¿Cómo diseñar y evaluar una arquitectura NoSQL multiinquilino que preserve una fuente
canónica comprensible, permita optimizar consultas mediante proyecciones derivadas,
pueda reconstruir esas proyecciones de forma verificable y ofrezca datos trazables para
minería de datos e inteligencia artificial?
\end{quote}

\section{Motivación}

\subsection{Motivación técnica}

El rendimiento de una aplicación conectada a una base remota no depende solamente del
tiempo de CPU. Importan también el número de solicitudes, los bytes transferidos, la
cantidad de registros procesados por el cliente y la variabilidad de la red. Una
proyección pequeña puede reducir el conjunto candidato de $N$ entidades a $K$
resultados relevantes, donde $K \ll N$. No obstante, esta mejora de lectura incrementa
el trabajo de escritura y almacenamiento. COBRIA busca hacer visible ese intercambio,
en vez de tratar cada índice como una optimización gratuita.

\subsection{Motivación de mantenibilidad}

En sistemas extensos, una entidad puede estar representada por un modelo de dominio, un
serializador, un repositorio, un servicio, reglas de validación, índices, catálogos y
documentación. Cuando esos artefactos evolucionan de manera independiente, aparecen
errores difíciles de detectar. Un enfoque dirigido por metadatos puede reducir la
duplicación de decisiones, siempre que no convierta los archivos generados en una falsa
medida de funcionalidad. COBRIA propone evaluar la generación por su capacidad para
preservar contratos y reducir defectos, no por la cantidad de código producido.

\subsection{Motivación analítica}

La minería de datos requiere transformar eventos y entidades en series, agregaciones,
secuencias, grafos o vectores de características. Cuando esas transformaciones se
consideran proyecciones versionadas, pueden reconstruirse desde datos canónicos,
compararse entre versiones y auditarse. Este enfoque conecta la arquitectura operacional
con prácticas de Almacén de características y MLOps, pero evita asumir que cualquier proyección
operacional es automáticamente adecuada para entrenar modelos.

\subsection{Motivación social y regional}

\ifnum\blindreview=1
El nombre COBRIA incorpora simbólicamente el origen regional de la investigación,
omitido durante la revisión ciega:
\else
El nombre COBRIA incorpora de manera simbólica la identidad de Coto Brus, Costa Rica:
\fi
\textbf{CO}--\textbf{BR}--\textbf{IA}. La referencia territorial no cambia el carácter
general del artefacto; expresa que la investigación tecnológica puede originarse fuera
de los centros tradicionales y producir conocimiento susceptible de evaluación y
reutilización internacional. El estudio evitará, sin embargo, convertir esa identidad
en una afirmación técnica o en evidencia de validez.

\section{La propuesta COBRIA en síntesis}

COBRIA significa \emph{Capa Orientada a Bases Reconstruibles para Inteligencia
Analítica}. La palabra \emph{objeto} se utiliza en un
sentido amplio: una entidad semántica identificable, no necesariamente una clase ligada
a un lenguaje de programación específico.

La arquitectura propone cinco movimientos principales:

\begin{enumerate}[leftmargin=*]
  \item representar el estado autorizado mediante entidades canónicas;
  \item separar nombres semánticos de la representación física almacenada;
  \item encapsular persistencia, reglas y operaciones mediante repositorios, servicios
  y casos de uso;
  \item definir índices, resúmenes, relaciones y conjuntos de datos como proyecciones con contrato
  de reconstrucción;
  \item conservar el linaje necesario para que minería de datos e IA utilicen
  representaciones versionadas, autorizadas y temporalmente correctas.
\end{enumerate}

\begin{figure}[H]
\centering
\resizebox{0.94\textwidth}{!}{%
\begin{tikzpicture}[
  node distance=7mm and 12mm,
  box/.style={draw=black,rounded corners=1pt,fill=white,
              minimum width=5.6cm,minimum height=9mm,align=center},
  derived/.style={draw=black,rounded corners=1pt,fill=black!6,
                  minimum width=4.4cm,minimum height=9mm,align=center},
  arrow/.style={-{Latex[length=2.4mm]},semithick,black}
]
  \node[box] (ui) {Interfaces y consumidores};
  \node[box,below=of ui] (uc) {Casos de uso y autorización};
  \node[box,below=of uc] (svc) {Servicios, políticas y composición};
  \node[box,below=of svc] (repo) {Repositorios / Mapeadores de datos};
  \node[box,below=of repo,fill=black!10] (canon) {Entidades canónicas};
  \node[derived,below left=of canon] (oper) {Proyecciones operacionales};
  \node[derived,below right=of canon] (anal) {Proyecciones analíticas y de IA};
  \draw[arrow] (ui) -- (uc);
  \draw[arrow] (uc) -- (svc);
  \draw[arrow] (svc) -- (repo);
  \draw[arrow] (repo) -- (canon);
  \draw[arrow] (canon) -- node[left,font=\scriptsize]{deriva} (oper);
  \draw[arrow] (canon) -- node[right,font=\scriptsize]{deriva} (anal);
  \draw[-{Latex[length=2.1mm]},dashed,black] (oper) to[bend right=24]
    node[below,font=\scriptsize]{verificación y reconstrucción} (canon);
  \draw[-{Latex[length=2.1mm]},dashed,black] (anal) to[bend left=24]
    node[below,font=\scriptsize]{linaje y reconstrucción} (canon);
\end{tikzpicture}
}
\caption{Vista conceptual inicial de COBRIA. Las proyecciones optimizan consumidores
específicos, pero su autoridad procede de las entidades canónicas.}
\label{fig:cobria-general}
\end{figure}

La flecha de reconstrucción en la \cref{fig:cobria-general} no implica que una
proyección sobrescriba libremente la entidad canónica. Indica que la proyección puede
compararse con el resultado esperado de su función de derivación y regenerarse cuando
sea necesario. Las correcciones del estado de negocio deben realizarse por medio de un
caso de uso autorizado.

\section{Vacío de investigación}

COBRIA se apoya en conceptos conocidos. Repositorio y Mapeador de datos separan el dominio de
la persistencia; las vistas materializadas aceleran consultas; CQRS distingue modelos
de lectura y escritura; las arquitecturas multiinquilino estudian la consolidación y el
aislamiento; los Almacenes de características organizan características para aprendizaje automático;
y la ciencia del diseño ofrece un método para construir y evaluar artefactos
\parencite{hevner2004designscience,peffers2007dsrm}.

El vacío que motiva este trabajo no es la inexistencia de esas ideas, sino la falta de
una articulación y evaluación conjunta para aplicaciones NoSQL basadas en estructuras
JSON que:

\begin{itemize}[leftmargin=*]
  \item trate la reconstruibilidad como contrato de toda proyección, no solo como una
  tarea de mantenimiento ocasional;
  \item conecte el modelo semántico y la representación física compacta mediante un
  mapeo versionado y reversible;
  \item incorpore el organización en identidad, rutas, consultas, autorización y pruebas;
  \item incluya proyecciones operacionales y analíticas bajo una taxonomía común;
  \item mida simultáneamente beneficios de lectura y amplificaciones de escritura y
  almacenamiento;
  \item preserve linaje desde la entidad operacional hasta características,
  predicciones y retroalimentación humana;
  \item sea contrastada en más de un tipo de sistema y exponga también los casos en que
  su complejidad no se justifica.
\end{itemize}

Por tanto, la originalidad pretendida es \emph{composicional y evaluativa}. El estudio
no afirmará que COBRIA sustituye todas las alternativas documentales ni que su
complejidad sea preferible para cualquier problema.

\section{Síntesis ampliada de trabajos relacionados}

La procedencia de datos antecede a COBRIA: la literatura distingue el origen de un
resultado y las partes de la fuente responsables de él \parencite{buneman2001provenance}.
COBRIA adopta esa preocupación, pero la lleva al contrato operacional de cada proyección:
fuente, versión, transformación, alcance y procedimiento verificable de reconstrucción.
Por ello, el linaje no se presenta como una invención, sino como una condición para que
una copia optimizada pueda auditarse y regenerarse.

La evolución de esquemas NoSQL también es un problema conocido. La ausencia de un
esquema global puede trasladar heterogeneidad y migraciones al código de aplicación, y
se han propuesto lenguajes declarativos para administrar esos cambios sin modificar el
motor \parencite{scherzinger2013schema}. COBRIA coincide en explicitar versiones, pero
añade la separación entre semántica canónica, mapeo físico y derivados reconstruibles.

En analítica, la arquitectura lakehouse busca reunir almacenamiento abierto, gobierno,
inteligencia de negocios y aprendizaje automático \parencite{armbrust2021lakehouse}.
Su escala es la plataforma de datos; COBRIA opera en otra capa: contratos de dominio y
proyecciones que pueden alimentar una plataforma analítica sin convertirla en autoridad
del estado transaccional. Ambas ideas pueden coexistir.

Finalmente, los sistemas de datos autoajustables estudian el uso del historial de carga
para planificar cambios físicos \parencite{pavlo2017selfdriving}. Esto delimita una línea
futura para COBRIA: recomendar o retirar proyecciones con un presupuesto observado. La
presente evaluación no implementa un optimizador autónomo; mide primero los beneficios y
costos que una política de ese tipo necesitaría comparar.

En conjunto, la revisión sitúa la novedad posible en la integración y en sus contratos,
no en renombrar vistas materializadas, procedencia, CQRS, repositorios o evolución de
esquema. La contribución es refutable porque puede fallar si la reconstrucción no es
exacta, si el aislamiento se rompe o si la amplificación supera el ahorro de lectura.

\section{Objetivos}

\subsection{Objetivo general}

Diseñar, formalizar y evaluar COBRIA como una arquitectura dirigida por metadatos,
basada en entidades canónicas y proyecciones reconstruibles, para determinar su efecto
sobre la eficiencia, integridad, mantenibilidad, recuperabilidad, aislamiento
multiinquilino y preparación para minería de datos e inteligencia artificial en sistemas
NoSQL.

\subsection{Objetivos específicos}

\begin{enumerate}[label=\textbf{OE\arabic*.},leftmargin=*]
  \item Definir los conceptos, componentes, responsabilidades y reglas de dependencia
  de COBRIA con terminología independiente de lenguajes y proveedores.

  \item Formalizar entidades canónicas, mapeos semántico--físicos, proyecciones,
  contratos de reconstrucción, alcance multiinquilino y linaje analítico.

  \item Clasificar las proyecciones en críticas, operacionales, analíticas y efímeras,
  y asociar a cada clase requisitos de consistencia y recuperación.

  \item Analizar la complejidad temporal y espacial de lectura, escritura, hidratación,
  mantenimiento y reconstrucción.

  \item Cuantificar latencia, transferencia de datos y amplificación de lectura,
  escritura y almacenamiento frente a alternativas sin proyecciones explícitas.

  \item Evaluar la exactitud y el costo de reconstrucciones totales e incrementales,
  incluida la detección de referencias huérfanas y proyecciones faltantes.

  \item Medir el efecto de metadatos compartidos sobre la evolución del esquema, la
  consistencia entre artefactos y el esfuerzo de mantenimiento.

  \item Evaluar el aislamiento multiinquilino mediante rutas acotadas, reglas de acceso y
  pruebas negativas con datos sintéticos.

  \item Analizar la capacidad de COBRIA para producir conjuntos de datos, series temporales y
  características versionadas con linaje y corrección temporal.

  \item Examinar la integración segura de modelos y agentes de IA mediante casos de
  uso autorizados, evidencia, consentimiento, auditoría y retroalimentación humana.

  \item Validar la aplicabilidad del artefacto en dos casos de estudio anónimos de
  diferente escala y naturaleza operacional.

  \item Identificar límites, antipatrones y condiciones bajo las cuales COBRIA resultaría
  innecesaria o excesivamente compleja.
\end{enumerate}

\section{Preguntas de investigación}

\begin{table}[H]
\centering
\caption{Preguntas de investigación de COBRIA.}
\label{tab:preguntas}
\begin{tabularx}{\textwidth}{>{\bfseries}p{1.6cm}X}
\toprule
Código & Pregunta \\
\midrule
RQ1 & ¿Cómo afecta COBRIA la complejidad y el comportamiento observado de las
operaciones de lectura, escritura y reconstrucción? \\
RQ2 & ¿En qué medida las proyecciones reducen la latencia, los registros procesados y
los bytes transferidos en consultas frecuentes? \\
RQ3 & ¿Qué amplificación de escritura y almacenamiento introduce el mantenimiento de
proyecciones y cuándo se justifica ese costo? \\
RQ4 & ¿Con qué exactitud, tiempo y consumo de recursos pueden reconstruirse las
proyecciones a partir de entidades canónicas? \\
RQ5 & ¿Cómo influye un mapeo semántico--físico dirigido por metadatos en la evolución
del esquema y en la consistencia entre modelos, validadores, reglas e índices? \\
RQ6 & ¿En qué medida el ámbito obligatorio mejora el aislamiento y limita las consultas
en un sistema multiinquilino? \\
RQ7 & ¿Puede la misma arquitectura adaptarse a un sistema orientado a eventos y
multimedia y a una plataforma empresarial con un dominio mucho más extenso? \\
RQ8 & ¿En qué medida las proyecciones analíticas versionadas mejoran la
reproducibilidad, el linaje y la corrección temporal de conjuntos de datos y características? \\
RQ9 & ¿Qué mecanismos arquitectónicos son necesarios para que modelos y agentes de IA
utilicen datos autorizados, auditables y acompañados de evidencia? \\
RQ10 & ¿Cuáles son los límites prácticos y las señales de sobrearquitectura de COBRIA? \\
\bottomrule
\end{tabularx}
\end{table}

Las preguntas combinan propiedades técnicas y de ingeniería. RQ1--RQ4 se responderán
principalmente mediante mediciones; RQ5 y RQ6 combinarán análisis de artefactos con
pruebas; RQ7 utilizará comparación entre casos; RQ8 y RQ9 abordarán preparación y
gobernanza de IA; y RQ10 sintetizará resultados negativos y costos de adopción.

\section{Hipótesis}

\begin{longtable}{>{\bfseries}p{1.6cm}p{5.2cm}p{7.1cm}}
\caption{Hipótesis y evidencia prevista.}\label{tab:hipotesis}\\
\toprule
Código & Hipótesis & Indicadores principales \\
\midrule
\endfirsthead
\toprule
Código & Hipótesis & Indicadores principales \\
\midrule
\endhead
H1 & Las consultas apoyadas en proyecciones procesan una cantidad de datos más cercana
al número de resultados que al tamaño total de la colección. & Registros leídos, bytes,
latencia p50/p95/p99 y amplificación de lectura. \\
H2 & La reducción del costo de lectura se obtiene a cambio de mayor amplificación de
escritura y almacenamiento. & Escrituras físicas por operación lógica, tamaño de
canónicos y derivados, duración de escritura. \\
H3 & Un contrato explícito permite reconstruir proyecciones deterministas con alta
exactitud a partir del estado canónico. & Cobertura de proyección, referencias
huérfanas, igualdad esperada--observada y tiempo de reconstrucción. \\
H4 & La reconstrucción incremental resulta más eficiente que la total cuando el número
de entidades modificadas es pequeño respecto al conjunto canónico. & Tiempo, lecturas,
escrituras y memoria para $\Delta N$ frente a $N$. \\
H5 & La definición compartida de metadatos reduce discrepancias entre modelo lógico,
representación física, validación, reglas e índices. & Defectos detectados, artefactos
modificados manualmente, tiempo de evolución y cobertura de contrato. \\
H6 & La incorporación obligatoria del organización en rutas, consultas y autorización reduce
la superficie de exposición cruzada. & Pruebas negativas, consultas fuera de ámbito,
registros innecesarios recuperados y violaciones detectadas. \\
H7 & Las proyecciones analíticas versionadas mejoran la reproducibilidad y el linaje de
conjuntos de datos frente a exportaciones ad hoc. & Tiempo de preparación, repetibilidad,
porcentaje de campos con procedencia y diferencias entre ejecuciones. \\
H8 & COBRIA puede aplicarse a los dos casos anónimos sin cambiar sus principios
centrales, aunque requiera proyecciones y políticas diferentes. & Correspondencia de
componentes, adaptaciones necesarias y resultados por caso. \\
H9 & El beneficio neto disminuye cuando una entidad tiene pocos registros, pocas
consultas o demasiadas proyecciones respecto a su valor funcional. & Función de costo,
presupuesto de proyección y análisis de sensibilidad. \\
\bottomrule
\end{longtable}

Las hipótesis son deliberadamente refutables. H2 y H9 reconocen costos y evitan evaluar
la arquitectura únicamente con métricas favorables. La investigación considerará
valioso identificar escenarios donde COBRIA no produzca una mejora neta.

\section{Contribuciones esperadas}

\subsection{Contribución conceptual}

Una definición neutral de COBRIA y un vocabulario que distinga entidad canónica,
representación física, proyección, índice, resumen, conjunto de datos y artefacto efímero. Esta
distinción busca facilitar la discusión entre desarrollo operacional, ingeniería de
datos y aprendizaje automático.

\subsection{Contribución formal}

Un modelo de entidades, ámbitos, mapeos y proyecciones; propiedades de autoridad,
reconstruibilidad, idempotencia y linaje; y condiciones bajo las cuales una proyección
puede considerarse verificable.

\subsection{Contribución arquitectónica}

Una integración reproducible de modelos de dominio, Repositorio/Mapeador de datos, Servicio
Capa, casos de uso, metadatos, proyecciones reconstruibles y componentes de analítica e
IA. El aporte arquitectónico incluye reglas de dependencia y una clasificación de
proyecciones por consistencia.

\subsection{Contribución algorítmica}

Pseudocódigo y posteriormente implementaciones de referencia para sincronización,
reconstrucción total e incremental, detección de inconsistencias, migración, generación
de conjuntos de datos y verificación de linaje.

\subsection{Contribución empírica}

Un conjunto de experimentos sobre dos casos anónimos que mida tanto beneficios como
costos. Se publicarán configuraciones y datos sintéticos suficientes para permitir la
replicación sin revelar información privada.

\subsection{Contribución práctica}

Una guía de adopción, un modelo de madurez, un presupuesto de proyecciones y criterios
para decidir cuándo utilizar COBRIA, cuándo simplificarla y cuándo preferir otra
organización documental.

\section{Casos de estudio y neutralidad del análisis}

La investigación se apoya en dos implementaciones existentes que comparten una familia
de decisiones arquitectónicas, pero difieren en dominio, escala y comportamiento. Para
evitar que el estudio se convierta en documentación de productos particulares, se
adoptan descripciones neutrales.

\begin{table}[H]
\centering
\caption{Caracterización inicial y anónima de los casos.}
\label{tab:casos}
\begin{tabularx}{\textwidth}{>{\bfseries}p{3.2cm}XX}
\toprule
Dimensión & Sistema A & Sistema B \\
\midrule
Naturaleza & Orientado a eventos, relaciones contextuales y artefactos multimodales. &
Plataforma empresarial multiinquilino con dominio extenso y operaciones transaccionales.\\
Escala estructural & Decenas de entidades y servicios especializados. & Cientos de
entidades, catálogos, schemas y proyecciones.\\
Datos destacados & Eventos temporales, secuencias, texto, métricas y multimedia. &
Entidades operacionales, financieras, configurables, de auditoría e integración.\\
Procesamiento & Servicios servicios internos, colas, procesos de trabajo y generación de artefactos. &
Operaciones autorizadas, índices, procesos de trabajo, políticas y generación por metadatos.\\
Valor para el estudio & Adaptabilidad multimodal, temporal, analítica y de IA. &
Escalabilidad estructural, multiinquilinato, gobierno del esquema y reconstrucción.\\
\bottomrule
\end{tabularx}
\end{table}

Los conteos exactos se fijarán en el protocolo experimental mediante scripts
reproducibles y una fecha de corte. No se utilizará el tamaño del código como sustituto
de calidad o completitud. Cuando una capacidad solo esté modelada o documentada, pero no
verificada mediante pruebas, se declarará como tal.

\section{Alcance}

El estudio incluye:

\begin{itemize}[leftmargin=*]
  \item sistemas NoSQL cuyo modelo de aplicación pueda representarse mediante
  documentos, árboles JSON o estructuras clave--valor jerárquicas;
  \item entidades canónicas y representaciones derivadas para consulta;
  \item repositorios, servicios y casos de uso como fronteras de responsabilidad;
  \item aislamiento lógico multiinquilino;
  \item mantenimiento y reconstrucción de índices y proyecciones;
  \item metadatos para modelos, esquemas, reglas, validación y documentación;
  \item complejidad computacional y mediciones de rendimiento;
  \item proyecciones analíticas, series temporales, características y linaje;
  \item integración gobernada de modelos y agentes de IA;
  \item dos casos de estudio anonimizados y una implementación neutral de referencia.
\end{itemize}

\section{Delimitaciones}

El estudio no pretende:

\begin{itemize}[leftmargin=*]
  \item demostrar que un motor NoSQL específico es superior a todos los demás;
  \item reemplazar los patrones Repositorio, CQRS, vistas materializadas, Almacén de características
  o arquitectura por capas;
  \item evaluar todos los motores NoSQL disponibles;
  \item probar que una cantidad elevada de entidades o archivos implica mayor madurez;
  \item presentar una implementación de IA como evidencia suficiente de exactitud,
  equidad o seguridad;
  \item entrenar modelos con información personal o sensible;
  \item revelar nombres, reglas comerciales, credenciales o datos de los sistemas
  utilizados como casos;
  \item afirmar causalidad más allá de lo que permita el diseño experimental.
\end{itemize}

Los experimentos iniciales se realizarán en entornos controlados y con datos sintéticos.
Las diferencias entre emuladores y producción se tratarán como amenaza a la validez
externa. Las integraciones externas no se considerarán funcionales únicamente porque
existan modelos o contratos en el código.

\section{Unidad de análisis y niveles de observación}

La unidad principal de análisis es una \emph{operación de datos COBRIA}: una acción que
lee o modifica una entidad canónica y, cuando corresponde, consulta o mantiene sus
proyecciones. El estudio observará cuatro niveles:

\begin{enumerate}[leftmargin=*]
  \item \textbf{Nivel de entidad:} identidad, campos, versión, organización e invariantes.
  \item \textbf{Nivel de operación:} lecturas, escrituras, transacciones, reintentos y
  efectos sobre proyecciones.
  \item \textbf{Nivel de sistema:} latencia, transferencia, almacenamiento,
  concurrencia, aislamiento y recuperación.
  \item \textbf{Nivel analítico:} conjunto de datos, característica, modelo, predicción, evidencia
  y retroalimentación.
\end{enumerate}

Esta separación evita mezclar conclusiones. Por ejemplo, que un acceso a caché tenga
costo promedio constante no implica que la operación distribuida completa sea
instantánea; y que un conjunto de datos sea reconstruible no implica que sea adecuado o ético
para entrenar un modelo.

\section{Consideraciones éticas iniciales}

La anonimización de los sistemas no es suficiente si los datos conservan combinaciones
que permitan reidentificar personas. Por ello, los experimentos utilizarán datos
sintéticos y no exportarán información de salud, identidad, pagos, conversaciones,
ubicaciones precisas ni contenido multimedia privado. Las pruebas de aislamiento se
realizarán únicamente en entornos autorizados.

En los componentes de IA se aplicarán, desde el diseño, los principios de trazabilidad,
supervisión humana, control de acceso y evaluación de riesgo. Marcos oficiales como el
marco de gestión de riesgos de inteligencia artificial de NIST enfatizan que la gestión del riesgo debe acompañar
el ciclo de vida completo del sistema, no limitarse al modelo entrenado
\parencite{nist2023airmf}. COBRIA no considerará una respuesta generada como una fuente
canónica de negocio mientras no pase por el caso de uso y las confirmaciones requeridas.

\section{Organización del documento maestro}

Después de esta introducción, el documento se organizará en siete bloques:

\begin{enumerate}[leftmargin=*]
  \item \textbf{Fundamentos y literatura:} conceptos necesarios y comparación con
  trabajos relacionados.
  \item \textbf{Diseño de COBRIA:} componentes, modelo formal, algoritmos, consistencia
  y reconstrucción.
  \item \textbf{Rendimiento:} complejidad, amplificación, modelo de costos y
  optimización.
  \item \textbf{Analítica e IA:} minería, Almacén de características, multimodalidad, RAG, agentes y
  gobernanza.
  \item \textbf{Implementación y casos:} implementación neutral y caracterización
  anónima de los dos sistemas.
  \item \textbf{Evaluación:} protocolo, resultados, análisis estadístico, discusión y
  amenazas a la validez.
  \item \textbf{Síntesis:} guía de adopción, conclusiones y líneas futuras.
\end{enumerate}

La estructura seguirá la lógica de la ciencia del diseño: el problema conduce a un
artefacto; el artefacto se formaliza e implementa; y su utilidad se juzga mediante
evaluación explícita, no por afirmación del autor \parencite{hevner2004designscience,
peffers2007dsrm}.

\section{Comparación sistemática con arquitecturas relacionadas}

La revisión ampliada distingue coincidencia de mecanismos y coincidencia de objetivos.
CQRS separa modelos de lectura y escritura, pero no obliga por sí mismo a declarar
linaje, ámbito o reconstrucción. Las vistas materializadas anticipan consultas y pueden
regenerarse desde sus fuentes, aunque su gobierno suele quedar dentro del motor
\parencite{microsoft2025cqrs,microsoft2025materialized}.
El abastecimiento por eventos conserva una secuencia autoritativa capaz de rehidratar
estado y proyecciones, pero COBRIA permite que la autoridad sea una entidad actual
versionada y no exige un registro completo de eventos \parencite{fowler2005eventsourcing}.

Los almacenes de características gobiernan definiciones para aprendizaje automático;
los sistemas de linaje explican procedencia; y los entornos tipo \emph{lakehouse}
unifican operación analítica. COBRIA conecta esos mecanismos con la frontera de casos de
uso, sin pretender sustituirlos \parencite{orr2021featurestores,li2023featurestore,
buneman2001provenance,armbrust2021lakehouse}. COBRIA trata las semánticas particulares
de persistencia, aislamiento e índices como capacidades explícitas de cada adaptador
NoSQL, no como supuestos universales.

\begin{table}[H]
\centering
\caption{Comparación entre COBRIA y enfoques relacionados.}
\label{tab:comparacion-literatura-ampliada}
\small
\begin{tabularx}{\textwidth}{p{2.7cm}p{2.2cm}p{2.1cm}X}
\toprule
Enfoque & Autoridad & Reconstrucción & Gobierno y diferencia principal \\
\midrule
CQRS & Parcial & Parcial & Separa lectura y escritura; ámbito y linaje no son inherentes.\\
Vista materializada & Externa & Sí & Optimiza consultas; el ámbito depende del motor y del diseño.\\
Abastecimiento por eventos & Eventos & Sí & El registro de eventos es autoridad y aporta historia.\\
Almacén de características & Fuentes analíticas & Parcial & Gobierna variables de modelos y su procedencia.\\
Arquitectura tipo \emph{lakehouse} & Archivos o tablas & Parcial & Unifica almacenamiento y procesamiento analítico.\\
COBRIA & Declarada & Obligatoria & Exige ámbito, linaje y gobierno desde operación hasta IA.\\
\bottomrule
\end{tabularx}
\end{table}

La matriz no establece superioridad. Muestra que la contribución propuesta reside en
hacer simultáneamente verificables autoridad, ámbito, reconstrucción, costo y continuidad
hacia inteligencia artificial. Si una solución existente satisface esos contratos con
menor complejidad, COBRIA recomienda conservar la solución más sencilla.

\section{Síntesis de la Parte I}

Esta primera parte estableció el problema que COBRIA pretende investigar. Los sistemas
NoSQL multiinquilino necesitan optimizar consultas sin perder autoridad, consistencia ni
capacidad de recuperación. Las necesidades analíticas añaden el requisito de conservar
linaje, versiones y corrección temporal desde la operación hasta la característica y la
predicción. COBRIA se propone como una integración de ideas existentes alrededor de dos
conceptos organizadores: \emph{datos canónicos} y \emph{proyecciones reconstruibles}.

La investigación no presupone que dicha integración sea siempre beneficiosa. Las
preguntas e hipótesis obligan a medir reducción de lectura junto con amplificación de
escritura y almacenamiento, exactitud de reconstrucción junto con su costo, y
reproducibilidad analítica junto con complejidad de adopción. Esta posición crítica será
la base del marco teórico, el modelo formal y la evaluación posterior.
